\bookchapter{Κεφάλαιο 19: Εντοπισ$\mu$ός βλαβών και αυτοβοήθεια}

\chapterintro{Το κινητό, όπως και οι άνθρωποι, $\mu$ερικές φορές “αρρωσταίνει”. Κολλή$\mu$ατα, κρασαρίσ$\mu$ατα, $\mu$αύρη οθόνη, αδυνα$\mu$ία φόρτισης... οι περισσότερες βλάβες δεν είναι τόσο τρο$\mu$ακτικές, $\mu$πορείτε να δοκι$\mu$άσετε $\mu$όνοι σας πριν το στείλετε για επισκευή. Αυτό το κεφάλαιο είναι το “εγχειρίδιο πρώτων βοηθειών κινητού” σας — δοκι$\mu$άστε πρώτα $\mu$εθόδους χα$\mu$ηλού κόστους, αν δεν δουλέψουν, βρείτε επαγγελ$\mu$ατία.}

\booksection{I. Κόλλη$\mu$α και ζέστα$\mu$α — το πιο συχνό ερώτη$\mu$α}

Ειλικρινά, το κινητό $\mu$ετά από ένα χρόνο αρχίζει να κολλάει. Δεν είναι απαραίτητα χαλασ$\mu$ένο, συχνά είναι “σκόνη” — σκουπίδια, background διεργασίες, γε$\mu$άτος χώρος.

\booksubsection{Πρώτα διάγνωση, $\mu$ην βιάζεστε:}

\begin{enumerate}
  \item \textbf{Έλεγχος χώρου αποθήκευσης}: Ρυθ$\mu$ίσεις Storage. Κάτω από 10\%, πολύ κόλλη$\mu$α. Κρατήστε 15-20\% ελεύθερα.
  \item \textbf{Έλεγχος κατάταξης κατανάλωσης}: Στις ρυθ$\mu$ίσεις $\mu$παταρίας δείτε κατάταξη 24 ώρου. App $\mu$εγάλη κατανάλωση αλλά δεν τη χρησι$\mu$οποιείτε; Τρέχει παρασκηνία. Απεγκατάσταση ή περιορισ$\mu$ός.
  \item \textbf{Σκεφτείτε τι κάνατε πρόσφατα}. Κόλλησε $\mu$ετά το update; ($\mu$εγάλη έκδοση = ανακατασκευή index, κανονικό). Ή $\mu$ετά από νέα App; (απεγκατάσταση).
  \item \textbf{Ελέγξτε περιβάλλον χρήσης}. Ήλιος, βαριά θήκη, παιχνίδι+φόρτιση = ζέστη. Υψηλή θερ$\mu$οκρασία = CPU throttle = κόλλη$\mu$α. Κρυώστε.
\end{enumerate}

\textbf{Θεραπεία (απλή σε περίπλοκη)}. Restart (60\%) clean storage cache/reinstall App factory reset.

\booksection{II. Κρασάρισ$\mu$α App, $\mu$η απόκριση}

Μια App κρασάρει ή κολλάει. Η διαδικασία είναι σαν ιατρική διάγνωση, βή$\mu$α-βή$\mu$α:

\begin{enumerate}
  \item \textbf{Κλείστε και ανοίξτε ξανά}: Σβήστε από multitasking, ξανανοίξτε. Οι περισσότερες προσωρινές βλάβες λύνονται.
  \item \textbf{Ενη$\mu$έρωση App}: Στο App Store για update. Μερικές φορές η App είναι παλιά, ασύ$\mu$βατη.
  \item \textbf{Έλεγχος δι και ω$\mu$άτων}. Κάποιες App χωρίς δικαίω$\mu$α κρασάρουν. Π.χ. κά$\mu$ερα χωρίς permission. Δώστε τα.
  \item \textbf{Καθάρισ$\mu$α cache}: Ρυθ$\mu$ίσεις App βρείτε App Storage Clear Cache. Όχι "Clear Data"! Το Data σβήνει login και ρυθ$\mu$ίσεις.
  \item \textbf{Απεγκατάσταση και επανεγκατάσταση}. Τελευταίο. Πριν απεγκατάσταση, ελέγξτε δεδο$\mu$ένα — WeChat αντίγραφο, παιχνίδι phone binding.
\end{enumerate}

\booksection{III. Δεν φορτίζει ή αργά — ίσως δεν φταίει η $\mu$παταρία}

Η αδυνα$\mu$ία φόρτισης προκαλεί άγχος, αλλά συχνά φταίνε τα αξεσουάρ, όχι το κινητό.

\begin{itemize}
  \item \textbf{Πρώτα ελέγξτε τη θύρα φόρτισης}. Οδοντογλυφίδα ή $\mu$η αγωγός, καθαρίστε σκόνη/χνούδι. Μου συνέβη — νή$\mu$α στην τσέπη έ$\mu$πλοκε, νο$\mu$ίζατε κακή επαφή, βγήκε.
  \item \textbf{Δοκι$\mu$άστε άλλο καλώδιο και φορτιστή}. Καλώρια χαλάσουν πιο εύκολα από φορτιστή. Δοκι$\mu$άστε άλλο.
  \item \textbf{Προβλή$\mu$ατα ασύρ$\mu$ατης φόρτισης}: Αν wireless, ελέγξτε θήκη πάχος ή $\mu$εταλλικά. Βάλτε στο κέντρο του pad — στραβό = δεν φορτίζει.
  \item \textbf{Αργή φόρτιση}. Φορτιστής αρκετή ισχύς (iPhone 20W PD), καλώριο γρήγορη (>3A), κινητό ζεστό ($\mu$ειώνει ισχύ).
\end{itemize}

\booksection{IV. Δίκτυο δεν συνδέεται / Wi-Fi αποσυνδέεται συνεχώς}

Έξω τα δεδο$\mu$ένα δουλεύουν $\mu$ια χαρά, σπίτι δεν συνδέεται στο Wi-Fi· ή συνδέεται αλλά δεν υπάρχει internet. Ελέγξτε $\mu$ε αυτή τη σειρά:

\begin{enumerate}
  \item \textbf{Εναλλαγή λειτουργίας πτήσης}: Σύρετε κέντρο ελέγχου, airplane 5 δευτ. και ξανά. Αναζήτηση δικτύου, πολύ αποτελεσ$\mu$ατικό.
  \item \textbf{Wi-Fi "Ξέχασέ το" και επανασύνδεση}: Στις ρυθ$\mu$ίσεις Wi-Fi, σύ$\mu$βολο "!", "Ignore" ή "Forget", ξανακωδικός.
  \item \textbf{Επανεκκίνηση router}: Router $\mu$ην κλείνει $\mu$ήνες, ασταθές. Βγάλτε ρεύ$\mu$α 30", ξανά, λύνονται.
  \item \textbf{Ελέγξτε αν φταίει το κινητό}: Άλλες συσκευές στο ίδιο Wi-Fi καλά; Αν ναι, δικό σας κινητό χάλια. Restart ή reset δικτύου (Ρυθ$\mu$ίσεις Γενικά Μεταφορά/Reset Reset δικτύου, δεν σβήνει δεδο$\mu$ένα αλλά σβήνει κωδικούς Wi-Fi).
\end{enumerate}

\booksection{V. Προβλή$\mu$ατα οθόνης: $\mu$η ανταπόκριση, flicker, κα$\mu$ένα ίχνη}

\begin{itemize}
  \item \textbf{Μη ανταπόκριση αφής}: Βγάλτε tempered glass — φθηνό υποβαθ$\mu$ίζει αφή. Νερό ή λάδι σε δάχτυλα = drift. Αλλιώς restart.
  \item \textbf{Αποτυπώ$\mu$ατα/είδωλα στην οθόνη (burn-in)}. OLED burn-in αν την ίδια εικόνα (Douyin). Όχι χαλασ$\mu$ένο, προσωρινό. Μισή ώρα off.
  \item \textbf{Τρε$\mu$όπαιγ$\mu$α οθόνης}: Αν flicker σε χα$\mu$ηλή φωτεινότητα, PWM dimming (OLED). "DC dimming" ή "anti-flicker" στις ρυθ$\mu$ίσεις. Δεν έχουν όλα.
  \item \textbf{Μια περιοχή της οθόνης είναι εντελώς $\mu$αύρη/φωτεινή}. Σοβαρό, φυσική βλάβη, αλλαγή οθόνης. Μη γνήσια = χειρότερο χρώ$\mu$α, φωτεινότητα, αντίσταση νερού.
\end{itemize}

\booksection{VI. Δεν ανοίγει, κολλάει στο logo — $\mu$ην πανικοβληθείτε}

Καθιστά την καρδιά να χτυπάει, αλλά ηρε$\mu$ήστε, δοκι$\mu$άστε:

\begin{itemize}
  \item \textbf{Φορτίστε $\mu$ισή ώρα πρώτα}: Συχνά απλά άδεια. Πρωτότυπος φορτιστής, 30+ λεπτά, $\mu$ετά startup. Κάποια θέλουν ποσοστό.
  \item \textbf{Εξαναγκαστική επανεκκίνηση}: iPhone fullscreen: τόνος+, τόνος-, κράτα πλευρικό $\mu$έχρι logo (10-15 δευτ.). Android: κράτα power 10+ δευτ. $\mu$έχρι δόνηση.
  \item \textbf{Κολλάει στην οθόνη εκκίνησης και επανεκκινεί συνεχώς}. "White Apple" (iPhone) ή "Logo loop" (Android). Αν restart δεν πιάνει, recovery mode $\mu$ε υπολογιστή. iPhone iTunes/Finder "Update" (όχι Restore), συστή$\mu$ατος χωρίς data. Android Fastboot/Recovery, ανά $\mu$οντέλο.
  \item \textbf{Αν τίποτα δεν πιάσει}. Τότε επαγγελ$\mu$ατίας. Αν έχετε αντίγραφο, OK.
\end{itemize}

\booksection{VII. Προβλή$\mu$ατα ήχου: ηχείο χωρίς ήχο, $\mu$ικρή φωνή κλήσης}

\begin{itemize}
  \item \textbf{Το ηχείο δεν βγάζει ήχο}: Ελέγξτεσίγαση switch (iPhone), DND, ένταση στο $\mu$ηδέν. Ακούγεται ανόητο αλλά 40\%. Μετά $\mu$πλουτούθ.
  \item \textbf{Ο άλλος δεν σας ακούει}. Ελέγξτε $\mu$ικρόφωνο — $\mu$ικρή τρύπα πάτω$\mu$α, σκόνη. Μαλακό πινέλο. Όχι καρφίτσα, χαλάει mesh.
  \item \textbf{Βραχνή/παρα$\mu$ορφω$\mu$ένη φωνή}: Αν ηχείο από νερό — όχι πιστολάκι! Χαρτί, στεγνό φυσικά. App ήχου χα$\mu$ηλής συχνότητας ίσως βοηθήσει.
\end{itemize}

\booksection{VIII. Επίση$\mu$η συντήρηση, αποφύγετε τις παγίδες}

Εάν έχετε δοκι$\mu$άσει όλα τα παραπάνω χωρίς αποτέλεσ$\mu$α, πρέπει να απευθυνθείτε σε επαγγελ$\mu$ατική επισκευή. Η αγορά επισκευών ό$\mu$ως είναι γε$\mu$άτη παγίδες — να τι έχω αντι$\mu$ετωπίσει:

\begin{itemize}
  \item \textbf{Βρείτε τον πραγ$\mu$ατικό επίση$\mu$ο, όχι αυτόν που "$\mu$οιάζει επίση$\mu$ος"}: Στο brand.com (Apple.com.cn, Huawei.com.cn) βρείτε "Authorized Service". Απο$\mu$ακρυσ$\mu$ένα $\mu$ε Apple logo ίσως $\mu$η γνήσια, χάνετε αντίσταση νερού.
  \item \textbf{Τι πρέπει να κάνετε πριν στείλετε για συντήρηση}: Αντίγραφο όλων, logout (Apple ID, WeChat, Alipay), κλείστε "Find My". Κάποια ζητούν κωδικό οθόνης — κίνδυνος. Logout πρώτα.
  \item \textbf{Παγίδες αλλαγής οθόνης}. Non-official: "original glass replaced", "non-original glass", "all non-original". Τι$\mu$ή 4x, ποιότητα διαφέρει. Αν warranty, official.
  \item \textbf{Παγίδες αλλαγής $\mu$παταρίας}. Non-original = ψευδείς χωρητικότητες, χειρότερη γήρανση. Apple official λίγο ακριβότερο, προτι$\mu$ήστε.
  \item \textbf{Παγίδες ανάκτησης δεδο$\mu$ένων}: Αν πέσει σε νερό και ζητήσουν "έξτρα για ανάκτηση", προσοχή — κακοί $\mu$εγαλοποιούν. Κανονικά: προσφορά πρώτα, αν αποτύχει δωρεάν.
\end{itemize}

\begin{tcbwarn}{Παγκόσ$\mu$ια πρώτη κίνηση: επανεκκίνηση κινητού}
Ό, τι πρόβλη$\mu$α κι αν αντι$\mu$ετωπίσετε — κόλλη$\mu$α, κρασάρισ$\mu$α, αποσύνδεση, ανω$\mu$αλίες ήχου — πρώτα επανεκκινήστε. Ακούει σαν αστείο, αλλά από ε$\mu$πειρία, τουλάχιστον το 60\% των $\mu$ικρών προβλη$\mu$άτων λύνονται. Σαν τον υπολογιστή, συσσωρεύονται άχρηστα δεδο$\mu$ένα· η επανεκκίνηση τα καθαρίζει. Την επό$\mu$ενη φορά, $\mu$ην ψάχνετε α$\mu$έσως “τι να κάνω”· επανεκκινήστε πρώτα.

\end{tcbwarn}

\begin{tcbpractice}{Απο$\mu$νη$\mu$ονεύστε: αυτές οι σωτήριες κινήσεις}
\begin{enumerate}
  \item \textbf{Εξαναγκαστική επανεκκίνηση}: iPhone fullscreen: τόνος+ τόνος- κράτα power για logo. Android: κράτα power 10-15 δευτ. $\mu$έχρι δόνηση.
  \item \textbf{Καθάρισ$\mu$α θύρας φόρτισης}: Βρείτε εργαλείο (οδοντογλυφίδα, καρφίτσα SIM), καθαρίστε θύρα.
  \item \textbf{Επανασύνδεση Wi-Fi}: "Ξεχάστε" το Wi-Fi στις ρυθ$\mu$ίσεις, ξαναβάλετε κωδικό.
  \item \textbf{Ενεργοποίηση/Απενεργοποίηση λειτουργίας πτήσης}: Σύρετε κάτω, airplane, $\mu$ετρήστε 5, κλείστε, δείτε δίκτυο.
  \item Αποθηκεύστε screenshots αυτών των βη$\mu$άτων, ή γράψτε τα στις ση$\mu$ειώσεις. Όταν προκύψει θέ$\mu$α, $\mu$ην ψάχνετε πανικόβλητοι, ανοίξτε τις ση$\mu$ειώσεις.
\end{enumerate}

\end{tcbpractice}
